\chapter{Introducción}
\label{chap:intro}

\section{Contexto y motivación}

Muchos mecanismos de elevación no presentan una carga constante al actuador. En una barra
articulada que se levanta tirando de un cable, la tensión que hace falta para sostener la carga
depende del ángulo que forman el cable y la barra; al girar la barra, esa geometría cambia y con
ella la tensión. Si el cable se enrolla en un tambor cilíndrico, el par que el motor debe
aportar, producto de la tensión por el radio, sigue fielmente esas variaciones. El actuador se
tiene que dimensionar entonces para el peor punto del recorrido, aunque la mayor parte del tiempo
trabaje muy por debajo de él.

Una solución clásica es cambiar la relación de transmisión a lo largo del recorrido haciendo que
el radio del tambor no sea constante. Es la idea del \emph{fusee} de la relojería del siglo XV,
un cono con una ranura helicoidal que compensaba la pérdida de par del muelle real a medida que se
desenrollaba \cite{landes1983}. La misma idea reaparece hoy en los tambores de radio variable de
la robótica de cables \cite{seriani2016,bieber2023} y en los sistemas pasivos de compensación de
gravedad \cite{shen2022,endo2010}. En este trabajo, al tambor en espiral se le llama
\emph{caracola} por su forma.

El proyecto en que se enmarca este TFG consiste en una caracola, accionada por un servomotor paso
a paso de lazo cerrado, que enrolla un cable y levanta una barra articulada con una masa de
\SI{2}{\kilo\gram}. Lo que se quiere demostrar es que el perfil de radio variable de la caracola
aplana el par que ve el motor frente a la variación de la tensión en el cable. Demostrarlo exige
dos cosas: un modelo que diga qué cabe esperar y un sistema de medida capaz de registrar a la vez
la posición del eje, la tensión del cable y el par del motor en ensayos repetibles. Construir ese
sistema ha sido la mayor parte del trabajo.

\section{Objetivos}

\subsection{Objetivo general}

Diseñar, construir y validar un sistema electrónico de accionamiento, control y adquisición para
un mecanismo de elevación con caracola de radio variable, y utilizarlo para comprobar
experimentalmente que la caracola reduce la variación del par motor frente a la variación de la
carga en el cable.

\subsection{Objetivos específicos}

\begin{enumerate}[label=\textbf{O\arabic*.}]
  \item Desarrollar un modelo analítico del mecanismo que relacione la geometría, la tensión del
        cable y el par en el tambor. A partir de él, deducir el perfil de radio que hace el par
        constante y compararlo con un tambor cilíndrico y con una espiral lineal.
  \item Diseñar una placa de circuito impreso que adapte el microcontrolador al entorno de la
        máquina: bus RS-485, entradas digitales protegidas para los finales de carrera, entradas
        analógicas acondicionadas, salidas de potencia y conexión de células de carga.
  \item Desarrollar el firmware del microcontrolador, capaz de gobernar el servo por Modbus RTU,
        leer la célula de carga, ejecutar el \emph{homing} y los ciclos automáticos, y vigilar los
        límites de seguridad con una temporización determinista.
  \item Definir un protocolo de comunicación robusto entre el microcontrolador y el ordenador,
        que transporte la telemetría por lotes y las órdenes con detección de errores.
  \item Desarrollar una aplicación de supervisión web para operar la máquina, visualizar los datos
        en tiempo real, grabar cada ensayo en un fichero reproducible y compararlo con el modelo.
  \item Caracterizar el sistema (temporización, repetibilidad de posición, comportamiento de la
        medida) y analizar los ensayos para cuantificar el efecto compensador de la caracola.
  \item Documentar los problemas técnicos encontrados y las soluciones adoptadas, así como la
        planificación, el presupuesto y el impacto del trabajo.
\end{enumerate}

\section{Alcance y limitaciones}

Entra en el alcance del trabajo todo el sistema electrónico y de software: la placa adaptadora, el
firmware, el protocolo, la aplicación de supervisión y el análisis de datos, además del modelo del
mecanismo y los ensayos. El diseño mecánico de la caracola y de la estructura se toma como dado
\completar{indicar si la caracola y la estructura se han diseñado y fabricado dentro de este TFG o
proceden de un trabajo previo}.

El trabajo tiene las siguientes limitaciones conocidas, que se discuten en los capítulos
correspondientes:
\begin{itemize}
  \item La posición del eje no se lee del codificador del accionamiento, sino que se obtiene
        integrando la velocidad que devuelve el propio servo (sección~\ref{sec:fw-posicion}).
  \item La célula de carga dispone sólo de una calibración provisional; sus valores en newtons se
        presentan como orientativos (sección~\ref{sec:res-limitaciones}).
  \item El par se obtiene en las unidades internas del accionamiento, cuya escala física no se ha
        confirmado. Por eso los resultados se expresan en variaciones relativas, que no dependen de
        esa escala.
  \item Los finales de carrera y el límite de fuerza se implementan por software; no hay una
        cadena de seguridad cableada (sección~\ref{sec:arq-seguridad}).
\end{itemize}

\section{Metodología}

El desarrollo ha sido incremental y guiado por pruebas. Cada subsistema (bus RS-485, célula de
carga, finales de carrera, entradas y salidas de la placa) se ha validado primero con un programa
de prueba aislado, que se conserva en el repositorio (anexo~\ref{anx:repo}), antes de integrarlo
en el firmware principal. El código se ha versionado con git desde el primer día; el historial de
cambios ha servido también para reconstruir la planificación real del proyecto
(capítulo~\ref{chap:gestion}). Los datos de todos los ensayos se guardan en crudo, con las
constantes de conversión en la cabecera de cada fichero, de modo que cualquier resultado se puede
recalcular sin repetir el ensayo. Las figuras de resultados de esta memoria se generan con un
programa reproducible (\fichero{memoria/figuras/generar\_figuras.py}) directamente a partir de esos
ficheros.

\section{Estructura del documento}

\begin{description}[style=nextline,leftmargin=2.2em]
  \item[Capítulo~\ref{chap:estado}] repasa los mecanismos de radio variable y las tecnologías de
        accionamiento, medida de fuerza, comunicación y supervisión relevantes.
  \item[Capítulo~\ref{chap:modelo}] desarrolla el modelo del mecanismo, el perfil ideal de par
        constante y el método experimental para separar las componentes del par.
  \item[Capítulo~\ref{chap:arquitectura}] recoge los requisitos y la arquitectura del sistema.
  \item[Capítulos~\ref{chap:hardware} a~\ref{chap:software}] describen, respectivamente, el hardware,
        el firmware, el protocolo de comunicación y la aplicación de supervisión.
  \item[Capítulo~\ref{chap:problemas}] reúne los problemas técnicos más relevantes y cómo se
        resolvieron.
  \item[Capítulo~\ref{chap:resultados}] presenta la caracterización del sistema y los resultados
        experimentales sobre la compensación de par.
  \item[Capítulo~\ref{chap:gestion}] trata la planificación, el presupuesto y el impacto.
  \item[Capítulo~\ref{chap:conclusiones}] expone las conclusiones y las líneas de trabajo futuro.
\end{description}

Los anexos contienen la asignación de pines, la especificación completa del protocolo, los
registros Modbus utilizados, el formato de los ficheros de ensayo, un manual de uso y la
estructura del repositorio.
